\begin{solapartat}
\punts{0,75} La força magnètica s'expressa com:
\[ \vec{F}_m = q\vec{v}\times\vec{B} \]

\figura[0.5]{sol_fig1.png}

Segons el dibuix podem comprovar que el camp magnètic i la velocitat són perpendiculars; el camp magnètic és perpendicular al pla horitzontal mentre que la velocitat està continguda al pla horitzontal, per tant, el mòdul del producte vectorial és:
\[ F_m = qvB \]

\figura[0.35]{sol_fig2.png}

D'altra banda, $\vec{F}_m$ és el resultat del producte vectorial, i per tant és perpendicular a $\vec{B}$ i a $\vec{v}$. Com que és perpendicular a $\vec{B}$, $\vec{F}_m$ es troba al pla horitzontal i com que és perpendicular a $\vec{v}$ llavors és perpendicular a la trajectòria, per aquesta raó les partícules dins els elèctrodes descriuen un moviment circular uniforme, MCU.

A més, segons la regla de la mà dreta i tenint en compte que la càrrega és positiva, la força magnètica va dirigida cap al centre dels semicercles.

Per tant, la força magnètica és igual a la massa per l'acceleració normal:
\[ F_m = qvB = m\frac{v^2}{r} \]

I, finalment, si agrupem els termes:
\[ \frac{v}{rB} = \frac{q}{m} \Rightarrow v = rB\frac{q}{m} \]

Si en la deducció no es justifica a partir de les direccions dels vectors que el producte vectorial dona $F_m = qvB$, cal restar 0,25 punts.

Si tampoc es justifica que la força magnètica és normal a la trajectòria, cal restar 0,25 punts.

\punts{0,25} Dins d'una D recorre una distància igual a mitja circumferència de radi $r$: $\pi r$. I com que es mou a una velocitat constant $v = rB\dfrac{q}{m}$, el temps que triga a passar per una D és:
\[ t = \frac{\pi r}{v} \]

D'altra banda, abans hem trobat:
\[ v = rB\frac{q}{m} \]

I substituint obtenim:
\[ t = \pi\frac{m}{Bq} \]

De la darrera expressió podem comprovar que el temps només depèn de $m$, $B$ i $q$, no depèn de la velocitat de la partícula.

\punts{0,15} Per tal que en l'espai entre elèctrodes la partícula s'acceleri, cal que el camp elèctric sigui paral·lel a la velocitat. Per tant, el camp elèctric ha de canviar la polaritat, tal i com s'indica a la figura de l'enunciat. És a dir, cal que el camp elèctric sigui altern.

\punts{0,1} A cada cicle, la partícula passa per les dues Ds; per tant, el temps que triga la partícula en passar per un elèctrode correspon a mig període del cicle
\[ T = 2\pi\frac{m}{Bq} \Rightarrow f = \frac{1}{T} = \frac{1}{2\pi}\frac{qB}{m} \]
\end{solapartat}

\begin{solapartat}
\punts{0,45} Abans hem trobat:
\[ v = rB\frac{q}{m} \]

Com que el protó surt quan arriba a l'extrem de l'elèctrode:
\[ v = RB\frac{q}{m} = 0,5\times 0,2\,\frac{1,602\times 10^{-19}}{1,67\times 10^{-27}} = 9,59\times 10^{6}\ \mathrm{m/s} \]

\punts{0,4}
\[ \lambda = \frac{h}{mv} = \frac{6,63\times 10^{-34}}{1,67\times 10^{-27}\times 9,59\times 10^{6}} = 4,14\times 10^{-14}\ \mathrm{m} \]

\punts{0,4} La velocitat és $0,1\times c = 3,00\times 10^{7}\ \mathrm{m/s}$
\[ R = \frac{v}{B}\frac{m}{q} = \frac{3,00\times 10^{7}}{0,2}\,\frac{1,67\times 10^{-27}}{1,602\times 10^{-19}} = 1,56\ \mathrm{m} \]
\end{solapartat}
